\section{The noise term and the environment index}
\label{app:eps}

This appendix records how the equations of \cite{ng2025} are read in this note.

\subsection{The index \texorpdfstring{$u$}{u}}
$u$ indexes environments; $Z^u=(Z_1^u,\dots,Z_n^u)$ is the latent vector of an environment-$u$ sample, so
the environment is marked on the random variable itself, exactly as it already is on $f_i^{(u)}$ and
$\epsilon_i^{(u)}$ in \cref{eq:anm} below. $\G_Z$ and $g$ carry no such superscript: the same DAG
and the same mixing function are used to generate $X^u=g(Z^u)$ in every environment. Write $p^{(u)}$ for
the density of $Z^u$, i.e.\ $p^{(u)}(z)=\Pr[Z^u\in dz]/dz$; this is the object the sufficient-change
assumptions are stated on, and it changes with $u$ precisely because the mechanisms $f_i^{(u)}$ and
$\epsilon_i^{(u)}$ do. No prior $p(u)$ over environments, and no mixture $\sum_u p(u)p^{(u)}(z)$, is
needed for \cref{thm:ng}: the theorem conditions on $u$ throughout and never marginalizes over it. (For
M1, a prior over which environment a given image comes from is supplied separately, by
\labelcref{item:P4}, and is used only for prediction, \cref{sec:use}; it is not part of the
identifiability argument.) In Fig.~1 of \cite{ng2025}, $u$ is drawn as a node with arrows into the
latents; we do not follow that convention here (\cref{fig:m1} draws $u$ as a plate index instead), since
an arrow $u\to Z_i$ would make $u$ itself a random variable with parents and children, which is neither
assumed nor needed.

In each environment, $p^{(u)}(z)=\prod_k p^{(u)}(z_k\mid\PA(z_k))$. The conditional
$p^{(u)}(z_i\mid\PA(z_i))$, which is the mechanism, is not the marginal $p^{(u)}(z_i)$ of one latent;
they coincide only for a node without parents. The sufficient-change assumptions are stated on the
derivatives of $\log p^{(u)}(z')$, the joint, which changes with $u$ because the mechanisms change.

\subsection{What \texorpdfstring{$\epsilon_i^{(u)}$}{epsilon} is in the ANM}
Eq.~(4) of \cite{ng2025}, with the environment marked on $Z$ as in \cref{sec:ng}, reads
\begin{equation}
Z_i^u=f_i^{(u)}\bigl(\PA(Z_i^u;\G_Z)\bigr)+\epsilon_i^{(u)},\qquad \epsilon_i^{(u)}\ \text{Gaussian}.
\label{eq:anm}
\end{equation}
Section~4 of the paper says only that the noise is Gaussian. It is not the standard normal: the
noise $\epsilon_i^{(u)}\sim\Nn\bigl(\mu_i^{(u)},(\sigma_i^{(u)})^2\bigr)$ carries a mean and a variance, and both
may depend on $u$. The evidence in the paper is the following.
\begin{itemize}[leftmargin=1.5em]
\item In the proof of Lemma~1 (Appendix~C.1), the conditional density of a sink is written with a scale
      $\sigma_i^{(u)}$, and $(\sigma_i^{(u)})^2$ is said to be the variance of $\epsilon_i^{(u)}$. The second
      derivative of the log-density of a sink with respect to itself is $-1/(\sigma_i^{(u)})^2$.
\item In the estimator of Section~4.3, the prior of each latent has a per-environment mean $\hat\mu_i^{(u)}$
      and scale $\hat\sigma_i^{(u)}$ computed from $u$ by small networks, and the scale does not depend on the
      parents.
\item In Section~4.2 and Section~6, the variances of the noise (and, in the simulations, the means) are
      generated randomly for each environment.
\end{itemize}
For a node with parents, the mean of the noise is absorbed into $f_i^{(u)}$; for a root node, $f_i^{(u)}$ is a
constant and the mean of the noise is the location of the node. The variance $(\sigma_i^{(u)})^2$ depends on $u$
but not on the parents.

The following is my own derivation, not a statement in the paper. For a sink $Z_i$, the entry
$\partial^2\log p^{(u)}/\partial Z_i^2=-1/(\sigma_i^{(u)})^2$ in $\vect w$ depends on $u$ only through the
variance. If the variance were the same in all environments, this entry would be constant in $u$, so it would
be zero in every difference vector, and \cref{as:2} could not hold. Hence \emph{the noise variance of every sink
latent must vary across environments}. For an isolated style latent this is the statement that the scale must
change across the style cells and not only the mean (cf.\ \cref{sec:reqs}).

\subsection{What \texorpdfstring{$\epsilon_i^{(u)}$}{epsilon} is in the HNM}
Eq.~(5) of \cite{ng2025} reads
\begin{equation}
Z_i^u=f_i^{(u)}\bigl(\PA(Z_i^u;\G_Z)\bigr)+\sigma_i^{(u)}\bigl(\PA(Z_i^u;\G_Z)\bigr)\,\epsilon_i^{(u)},
\label{eq:hnm}
\end{equation}
with $\epsilon_i^{(u)}$ Gaussian and $f_i^{(u)}$, $\sigma_i^{(u)}$ third-order differentiable. Here the Gaussian
noise is the standard normal: in the proof of Lemma~2 (Appendix~C.2) the density is written in terms of
$\bigl(Z_i^u-f_i^{(u)}\bigr)/\sigma_i^{(u)}(\PA(Z_i^u))$, so the scale is carried by the function
$\sigma_i^{(u)}$ of the parents, and this function also depends on $u$.

\subsection{Summary of the difference}
\begin{center}
\begin{tabular}{@{}lll@{}}
\toprule
 & ANM, Eq.~(4) & HNM, Eq.~(5)\\
\midrule
noise & $\epsilon_i^{(u)}\sim\Nn(\mu_i^{(u)},(\sigma_i^{(u)})^2)$ & $\epsilon_i^{(u)}\sim\Nn(0,1)$\\
scale & $\sigma_i^{(u)}$ depends on $u$ only & $\sigma_i^{(u)}(\PA(Z_i))$ depends on $u$ and on the parents\\
sink $Z_i$ & $\partial^2\log p/\partial Z_i^2=-1/(\sigma_i^{(u)})^2$ & $\partial^2\log p/\partial Z_i^2=-1/\sigma_i^{(u)}(\PA(Z_i))^2$\\
sink property & $\partial^3\log p/\partial Z_i^2\partial Z_j=0$ for all $j$ & $\partial^3\log p/\partial Z_i^2\partial Z_j=0$ only for $Z_j\notin\PA(Z_i)$\\
\bottomrule
\end{tabular}
\end{center}

\subsection{Two common misreadings}
\begin{itemize}[leftmargin=1.5em]
\item \emph{The index $u$ does not mean that the noise is resampled per environment in a special way.} Within
      one environment every sample receives fresh, independent noise for every node; $u$ only selects the
      parameters of the distribution from which the noise is drawn. If the noise did not depend on $u$, it
      would still be drawn afresh for each sample and each node, independently across nodes. Sharing one draw
      among the nodes would break the independence of the noise terms.
\item \emph{``Gaussian'' does not mean standard normal in the ANM.} It does in the HNM, where the scale is a
      separate function.
\end{itemize}
In the simulations of \cite{ng2025} the means and variances of the noise and the weights of the mechanisms are
drawn once per environment, which generates a family of environments.

\subsection{M1 in these terms}
For a content latent in the cell $c=(y,t)$, each image receives a fresh draw from
$\Nn(0,(\sigma_{c,i}^{(c)})^2)$, and the same variance applies at every location; the mean is carried by
$f_i^{(c)}$. For a style latent in the cell $r=(t,e)$, each image receives a fresh draw from
$\Nn(\mu_j^{(r)},(\sigma_{s,j}^{(r)})^2)$.
